\chapter*{Conclusion Générale}
\markboth{Conclusion Générale}{}
\addcontentsline{toc}{chapter}{Conclusion Générale}

Ce projet a été consacré à la conception et au développement d'un système de détection
des textes générés par l'intelligence artificielle, dans un contexte où la distinction entre
productions humaines et contenus automatiques devient un enjeu croissant dans les
domaines académiques, professionnels et médiatiques. En suivant la méthodologie
CRISP-DM, nous avons mené chaque étape du projet de manière structurée : de la
compréhension de la problématique jusqu'au déploiement de l'application, en passant
par la collecte et la préparation d'un jeu de données de 487\,235 textes, la modélisation
et l'évaluation de deux architectures de deep learning. Le prototype développé démontre
la faisabilité technique de notre approche et illustre la pertinence d'un tel outil pour
détecter automatiquement l'origine d'un texte.
\\
\\
Les résultats obtenus confirment la supériorité du modèle CNN, qui atteint une accuracy
de 99{,}8\% contre 96\% pour le modèle LSTM, grâce à sa capacité à identifier des
motifs stylistiques et lexicaux caractéristiques des textes générés par IA. Ces performances
ont été rendues possibles par une préparation rigoureuse des données, notamment le
rééquilibrage des classes et la tokenisation via BERT. Toutefois, des limites subsistent :
la sensibilité du modèle à la qualité du jeu de données, la difficulté à généraliser à des
textes produits par des modèles génératifs plus récents, et la nécessité d'améliorer
l'interprétabilité des prédictions. Ces constats ouvrent la voie à des perspectives futures,
notamment l'adoption d'architectures Transformer pré-entraînées, l'enrichissement du
corpus d'entraînement et l'adaptation à différents contextes linguistiques.
\\
\\
En définitive, ce travail constitue une première étape vers la mise en place d'outils
fiables permettant de mieux encadrer l'usage de l'intelligence artificielle dans la
production écrite. L'application déployée sur Render, combinant un backend FastAPI
et un frontend React, offre une solution accessible et fonctionnelle. Il met en évidence
l'importance de continuer à explorer ce domaine en constante évolution, afin d'apporter
des solutions toujours plus robustes et adaptées aux défis posés par les technologies
émergentes.